% chapters/conclusions.tex
\addcontentsline{toc}{section}{ВИСНОВКИ}
{\centering\bfseries ВИСНОВКИ\par}
\vspace{1em}

У бакалаврській кваліфікаційній роботі розв'язано актуальну задачу
розроблення веб-симулятора 16-точкового ШПФ radix-4 з інтерактивною
візуалізацією потокового виконання, який поєднує достовірність обчислення
з наочним відображенням dataflow-природи алгоритму.

\textbf{Основні результати роботи.}

\begin{enumerate}
  \item Проаналізовано сімейство алгоритмів ШПФ та встановлено, що алгоритм
    radix-4 для $N = 16$ потребує лише 9~нетривіальних комплексних множень
    і 80~комплексних додавань проти 256~множень у прямому ДПФ; прискорення
    за числом нетривіальних операцій --- $\approx 28{,}4\times$
    (підрозд.~1.1, 4.3). Виконано системний аналіз існуючих рішень (FFTW,
    CMSIS-DSP, cuFFT, SPIRAL) і обґрунтовано вибір клієнтського SPA на
    React~Flow як засобу, що поєднує достовірність обчислення з доступністю
    без встановлення (підрозд.~1.2, 1.3).

  \item Розроблено повну математичну модель тристадійного алгоритму
    $4 \times 4$ декомпозиції: відображення часових і частотних індексів
    $n = 4s + r$, $k = 4k_2 + k_1$~(\ref{eq:idx_map}); три стадії
    обчислення~(\ref{eq:stage1})--(\ref{eq:stage3}); таблицю поворотних
    множників $W_{16}^{r k_1}$; оптимізований 4-точковий ДПФ без
    нетривіальних множень~(\ref{eq:dft4_butterfly}) (підрозд.~2.2).

  \item Розроблено граф потоку даних обчислювача (56~вузлів чотирьох
    типів \texttt{source}, \texttt{dft4}, \texttt{twiddle}, \texttt{sink}
    та 64~дуги з \texttt{delay}~$=600$). Визначено критичний шлях
    $T_{\mathrm{kr}} = 3400$~умовних одиниць~(\ref{eq:critical_full}) і
    теоретичний коефіцієнт паралелізму $K_p^{\text{теор}} \approx 2{,}9\times$~(\ref{eq:parallelism})
    завдяки одночасній активації чотирьох незалежних \texttt{dft4}-вузлів
    на кожному ярусі (підрозд.~2.3).

  \item Реалізовано веб-симулятор на стеку Next.js~14 / TypeScript~5 /
    React~Flow~12 / Zustand~4 з ядром \texttt{DataflowRuntime}, що реалізує
    правило активації токенів через чотирикроковий метод
    \texttt{tick()}~(\ref{eq:delivery})--(\ref{eq:emit_time}).
    Функція \texttt{generateDFT4x4()} програмно будує граф відповідно до
    математичної моделі розділу~2. Анімація токенів реалізована через
    інтерполяцію прогресу~(\ref{eq:progress})--(\ref{eq:lerp}), теплова
    карта активності --- через затухання~(\ref{eq:heat_decay}) з коефіцієнтом
    $\alpha = 0{,}92$ за кадр (розд.~3).

  \item Верифіковано коректність симулятора на шести тестових
    сценаріях~(таблиця~\ref{tab:scenarios}): функція \texttt{compareWithFFT()}
    не виявила жодного розбіжного спектрального відліку; максимальна
    абсолютна похибка не перевищила $3 \times 10^{-12}$, що на три порядки
    менше від порогу прийняття $\varepsilon = 10^{-9}$. Підтверджено
    ермітову симетрію $X(N-k) = X^*(k)$ для дійсних вхідних послідовностей
    (підрозд.~4.3).

  \item Експериментально підтверджено паралелізм dataflow-моделі: практичний
    коефіцієнт $K_p \approx 1{,}88\times$ з урахуванням затримок
    дуг~(\ref{eq:parallelism_full}) і теоретичний
    $K_p^{\text{теор}} \approx 2{,}9\times$ без них. Метрика
    \texttt{latencyEmaMs} у стаціонарному режимі точно збігається з
    теоретичним критичним шляхом~(\ref{eq:critical_full}) (підрозд.~4.4).

  \item Усі вісім нефункціональних вимог NF1--NF8 (інтерактивність
    $\geq 60$~fps, складність $O(N \log N)$, пам'ять $< 100$~МБ, точність
    $< 10^{-9}$, кросбраузерність, MIT-ліцензія, портативність, тестове
    покриття $> 80\%$) підтверджено експериментально
    (таблиця~\ref{tab:nf_verify}) у трьох браузерах
    (Chrome, Edge, Firefox).
\end{enumerate}

\textbf{Практичне значення.}
Розроблений симулятор є навчально-дослідницьким інструментом для вивчення
потокових алгоритмів цифрового оброблення сигналів у курсах «Цифрове
оброблення сигналів», «Архітектура обчислювальних систем» та «Паралельні
обчислення» на бакалаврській програмі 122 «Комп'ютерні науки». Застосунок
виконує реальний алгоритм radix-4 (а не його імітацію), доступний через
будь-який сучасний веб-браузер без встановлення додаткового програмного
забезпечення та підтримує регульовану швидкість, покроковий режим і
пресети вхідних сигналів, що робить його зручним для демонстрації під час
лекційних занять. Архітектура DFG (56~вузлів, 64~дуги, параметри
\texttt{latency} та \texttt{delay}) може слугувати вихідною специфікацією
для апаратної реалізації обчислювача ШПФ на FPGA, де граф потоку даних
безпосередньо відображається на маршрутизацію сигналів між DSP-блоками.
Відкритий код проекту (MIT) дозволяє його адаптацію і розширення іншими
дослідницькими групами.

\textbf{Обмеження дослідження.}
Симулятор обмежений фіксованим розміром перетворення $N=16$; систематичне
тестування у браузері Safari та на мобільних пристроях не проведено;
довготривалий профайлінг пам'яті (понад 24~години безперервної роботи)
не виконувався; WebWorker-ізоляція обчислювального ядра від UI-потоку
не реалізована.

\textbf{Подальші напрями вдосконалення.}

\begin{enumerate}
  \item Масштабування моделі на $N \in \{64, 256, 1024\}$ через ієрархічне
    застосування \texttt{generateDFT4x4()} з рекурсивним вкладенням;
    оцінка меж продуктивності React Flow для великих графів
    ($>500$~вузлів) та впровадження віртуалізації рендерингу.
  \item Перенесення ядра \texttt{DataflowRuntime} у WebWorker для повного
    відокремлення обчислень від UI-потоку та забезпечення плавної
    анімації при великих $N$.
  \item Підтримка двовимірного ШПФ для задач оброблення зображень
    (медична візуалізація, спектральна фільтрація) шляхом послідовного
    застосування 1D-симулятора рядково та стовпчиково.
  \item Інтеграція з LMS-платформами університету через LTI-протокол
    для вбудовування симулятора в курси дистанційного навчання;
    серверне збереження пресетів і результатів експериментів.
  \item Розширення бібліотеки алгоритмів: підтримка radix-2 (для
    порівняння), mixed-radix, Bluestein, Winograd; порівняльна
    візуалізація різних стратегій декомпозиції на одному графі.
  \item Апаратна верифікація DFG-специфікації на платформі FPGA
    (Xilinx Artix-7 або Intel Cyclone V) з використанням
    \texttt{generateDFT4x4()} як джерела генерації RTL-опису.
\end{enumerate}
